\documentclass{article}
\usepackage[spanish]{babel}
\usepackage[utf8]{inputenc}
\usepackage[T1]{fontenc} % allows to copy accented words on PDF
\usepackage{blindtext}
\usepackage[letterpaper, margin=1in]{geometry}
\usepackage{url}
\usepackage{pgfplots}
\pgfplotsset{width=10cm,compat=1.9}
\usepackage{graphicx}
\usepackage{booktabs}
\usepackage{amsmath}
\usepackage{listings}
\usepackage{lmodern}
\usepackage[framed,numbered,captionpos=t]{matlab-prettifier}
\usepackage{color} %red, green, blue, yellow, cyan, magenta, black, white
\definecolor{gray97}{gray}{.97}
\definecolor{gray75}{gray}{.75}
\definecolor{gray45}{gray}{.45}
\definecolor{mygreen}{RGB}{28,172,0} % color values Red, Green, Blue
\definecolor{mylilas}{RGB}{170,55,241}
%----Java----
\definecolor{mygreen}{rgb}{0,0.6,0}
\definecolor{mygray}{rgb}{0.5,0.5,0.5}
\definecolor{mymauve}{rgb}{0.58,0,0.82}
\lstset{ %
  backgroundcolor=\color{white},   % choose the background color
  style      = Matlab-editor,
  basicstyle = \mlttfamily,
  numbers=left,
  stepnumber=1,
  xleftmargin=.1\textwidth,
  xrightmargin=.1\textwidth,
  %basicstyle=\footnotesize,        % size of fonts used for the code
  breaklines=true,                 % automatic line breaking only at whitespace
  captionpos=t,                    % sets the caption-position to bottom
  commentstyle=\color{mygreen},    % comment style
  escapeinside={\%*}{*)},          % if you want to add LaTeX within your code
  keywordstyle=\color{blue},       % keyword style
  stringstyle=\color{mymauve},     % string literal style
  morekeywords={run, getDimensions, newArray, getResult}
}

% -=-=-=-=-=-=-=-=-=-=
% Cabecera y pie de página
\usepackage{fancyhdr}
\pagestyle{fancy}
\fancyhf{}
\rhead{Práctica 1: Nombre de la práctica}
\lhead{Instituto Tecnológico de Morelia}
\cfoot{\thepage}
% -=-=-=-=-=-=-=-=-=-=


% -=-=-=-=-=-=-=-=-=-=
% Defniciones para título y autores
\makeatletter
\def\HUGE{\@setfontsize\HUGE{30}{60}}
\setlength{\parskip}{0.1in}

\newcommand{\myTitle}{Nombre de la una Práctica}

\newcommand{\authora}{Nombre del Alumno 1}
\newcommand{\authorb}{Nombre del Alumno 2}

\newcommand{\myReviewers}{Dr. Gerardo Marx Chávez-Campos}

\newcommand{\myDate}{Agosto del 2023}

\begin{document}

% archivo de la portada
\include{TitlePage}

% índice de contenido
\tableofcontents




\newpage

\section{Introducción}

% Considere que las referencias están en formato \textit{Bibitems} y se incluyen al final de este archivo.

Este es un ejemplo de como se deben insertar referencia en el documento~\cite{conagua-2011}. Considere que las referencias están en formato \textit{Bibitems} y se incluyen al final de este archivo~\cite{conagua-2011,undesa2008}.


% Texto dummy para rellenar el documento
\blindtext[2]

 % ejemplo de figura
\begin{figure}[htb]
    \centering
    \includegraphics[width=3.5in]{universe}
    \caption{Pie de texto en la figura.}
    \label{fig:universe}
\end{figure}

\blindtext[3]

\section{Objetivos}

\subsection{Objetivo general}
\blindtext[1]

\subsection{Objetivos específicos}

\begin{itemize}
	\item Primer elemento y su descripción relacionada con los objetivos.
	\item Segundo elemento y su descripción relacionada con los objetivos.
	\item Tercer elemento y su descripción relacionada con los objetivos.
	\item Cuarto elemento y su descripción relacionada con los objetivos.
\end{itemize}

\section{Marco Teórico}

\blindtext[2]

El Error Cuadrático medio queda definido por la Ecuación \eqref{eq:mse}:

\begin{equation}
	MSE(X,h_\theta)=\frac{1}{m}\sum_{i=0}^m\left(\theta x^{(i)}-y^{(i)} \right)^2 \label{eq:mse}
\end{equation}

\section{Desarrollo}

\blindtext[1]

\begin{lstlisting}[
  language=Python,
  label=lst:model,
  frame= lines,
  caption= {Script to implement the Joule heating ODE model considering physical geometry and its thermal and electrical properties}]
# # Model definition
# the main equation is:
# $$Q_g+Q_i=Q_a+Q_l

import numpy as np
from scipy.integrate import odeint
import termplotlib as tpl
import matplotlib.pyplot as plt

def jouleODE(y,t):
    rho = 7870 # kg/m^3, material's density
    Cp = 800 # J/(KgC)800, Heat Capacity
    rhoE = 1.15E-6 #OhmXm @850C
    h = 170 # W/m^2K
    delta_z = 0.026 #m
    diameter = 0.003 #m 
    Ac = 2*3.141592*(diameter/2)*(diameter/2)
    alpha = rhoE*delta_z/Ac
    V = Ac*delta_z
    beta = rho*Cp*V
    A_cylinder = 2*3.141592*diameter/2*delta_z
    A_s = 2*Ac+A_cylinder
    theta = A_s*h
    Tm = 27 #C room's temp
    I = 140 #Amps, input's current
    dydt = alpha/beta*(I*I)-theta/beta*(y-Tm)
    return dydt

y0 = 27
t = np.linspace(0,1500)
y = odeint(jouleODE, y0,t) 
plt.plot(t,y,'.:k', label='model')
plt.xlabel('Time (sec)')
plt.ylabel('y(t)')
plt.legend()
plt.show()
\end{lstlisting}



\section{Resultados}

\blindtext[1]

Plotting from data:

\begin{figure}[htb]
\centering
\begin{tikzpicture}
\begin{axis}
[
	title={Resultados de la medición},
	xlabel={$x$ Temperature [\textcelsius]},
  ylabel={Solubility [g per 100 g water]},
    xmin=0, xmax=100,
    ymin=0, ymax=120,
    xtick={0,20,40,60,80,100},
    ytick={0,20,40,60,80,100,120},
    legend pos=north west,
    ymajorgrids=true,
    grid style=dashed
]

\addplot[
    color=blue,
    mark=square,
    ]
    coordinates {
    (0,23.1)
    (10,27.5)
    (20,32)
    (30,37.8)
    (40,44.6)
    (60,61.8)
    (80,83.8)
    (100,114)
    };
    \legend{Data $f(x)$}
\end{axis}
\end{tikzpicture}
\end{figure}


\section{Conclusiones}

\blindtext[1]



\clearpage
%----------
\begin{thebibliography}{99}

\bibitem{conagua-2011}
Comisión Nacional del Agua.
\textit{Estadísticas del Agua en México, 2011.} Recuperado 8 de agosto de 2021, de \url{http://www.conagua.gob.mx/CONAGUA07/Publicaciones/Publicaciones/SGP-1-11-EAM2011.PDF}

\bibitem{undesa2008}
United Nations. Department of Economics and Social Affairs, Population Department,
\textit{World Population Prospects: The 2008 Revision: Medium fertility variant 2010 - 2050.}
Consultado en: \url{http://esa.un.org/unpd/wpp2008/index.html} 

\bibitem{fondo2021}
Fondo para la Comunicación y la Educación Ambiental, A.C.
\textit{Contaminación en México.} Recuperado 10 de agosto de 2021, de \url{https://agua.org.mx/agua-contaminacion-en-mexico/}

\bibitem{fernandez2020}
Olaiz-Fernández GA, Gómez-Peña EG, Juárez-Flores A, Vicuña-de Anda FJ, Morales-Ríos JE, Carrasco OF. \textit{Panorama histórico de la enfermedad diarreica aguda en México y el futuro de su prevención. Salud Publica Mex [Internet]}. 20 de diciembre de 2019 [citado 19 de agosto de 2021];62(1, ene-feb):25-3. Disponible en: \url{https://saludpublica.mx/index.php/spm/article/view/10002}

\bibitem{soto2016}
Soto-Estrada, Guadalupe; Moreno-Altamirano, Laura  y  Pahua-D\'iaz, Daniel. \textit{Panorama epidemiológico de México, principales causas de morbilidad y mortalidad.} Rev. Fac. Med. (Méx.) [online]. 2016, vol.59, n.6 [citado  2021-08-19], pp.8-22, ISSN 2448-4865. Disponible en: \url{http://www.scielo.org.mx/scielo.php?script=sci_arttext&pid=S0026-17422016000600008&lng=es&nrm=iso}. 

\bibitem{camara2021}
Cámara de Diputados. \textit{Constitución Politica de los Estados Unidos Mexicanos.} Artículo 4, consultado en Agosto del 2021 en:  
\url{http://www3.diputados.gob.mx/camara/001_diputados/012_comisioneslxii/01_ordinarias/002_agua_potable_y_saneamiento/13_marco_juridico/01_constitucion_politica_de_los_estados_unidos_mexicanos}

\bibitem{merida2020}
Mérida-Cano, Marvin Eduardo.
\textit{Calidad bacteriológica del agua y su relación con el potencial de óxido reducción (ORP)}, 2020.

\bibitem{astudillo2019}
Astudillo-Naderas, Ricardo y Montes de Oca Barrera, Fabiola. \textit{Electrodo de grafito como electrodo indicador de reacciones redox}.  Universidad Nacional Autonoma de México, Facultad De estudios Superiores Cuautitlán, 2019.

\bibitem{steininger1996}
Steininger, Jacques y Pareja, Catherine. 
\textit{Respuesta Del Sensor De ORP En Agua Clorada,} 1996

\bibitem{lin2017}
Wen-Chi, Lin; Klaus Brondum, Charles W. Monroe; y Mark A. Burns. \textit{Monitoreo del pH, potencial de oxidación-reducción (ORP) y conductividad de agua,}  2017

\bibitem{trevor2004}
Suslow, Trevor V. \textit{Potencial de oxidación-reducción (ORP) para el monitoreo, control y documentación de la desinfección del agua}, 2004.
\end{thebibliography}

\end{document}
